\begin{enumerate}
\setlength\itemsep{1 in}

\item 
We have gone through several kinds of equations now and lets sum up some of these \emph{as proportions}:
\begin{itemize}
	\setlength\itemsep{.5 in}
	\item acceleration is \rule{1in}{0.1mm} net force.
	\item assuming constant acceleration and beginning at rest, an object's velocity is \blank the displacement.
	\item assuming constant acceleration and beginning at rest, an object's displacement is \blank the elapsed time.
	\item for an object that has been dropped, the distance it has fallen is \blank its velocity at that distance.
\end{itemize} 
\medskip


\item
I push a \SI{100}{kg} box starting at rest along a friction-less floor, with a force of \SI{100}{\newton} for \SI{10}{\second}. How fast is the box going at this point? If I did the same thing to a \SI{200}{kg} box, then how fast is it going after \SI{10}{s}?
\hugeskip

\item
I push a \SI{100}{kg} box starting at rest along a friction-less floor, with a force of \SI{100}{\newton} over a distance of \SI{10}{m}. How fast is the box going at this point? If I did the same thing to a \SI{200}{kg} box, then how fast is it going after \SI{10}{m}?\\
Some starters:
\begin{itemize}
	\item What is the net force on the box?
	\item What is the acceleration of the box?
	\item What is the final velocity after \SI{10}{\meter}
\end{itemize}
\hugeskip


\item
Following up on the previous problem, if I stopped pushing after \SI{10}{m} and the box continued with its speed, and then at then started sliding up a \ang{20} ramp, then how far along the length of the ramp would the box rise? What height is this above the horizontal?  Do the \SI{100}{kg} and the \SI{200}{kg} box rise to the same height?\\
Some starters:
\begin{itemize}
	\item What is the net force on the box as it goes up the inclined plane?
	\item What is the acceleration of the box as it goes up the inclined plane?
	\item What is the sign of the displacement of the box going up the plane?
	\item Is the sign of acceleration the same or different than displacement? 
\end{itemize}
\hugeskip

\item A 50-kg crate of apples is lowered by a rope straight down and has an acceleration of \SI{1.0}{m/s^2} in the downward direction. What is the tension in the rope? State its magnitude and direction. \bigskip

\item A \SI{900}{kg} elevator moves downward \SI{11.0}{\meter} in \SI{4}{\second}. The tension in the supporting cable is \SI{7800}{\newton}. What are the initial and final velocities of the elevator?
\bigskip



%\item 
%We want to examine the idea of \emph{mechanical advantage}. 
%
%\begin{itemize}
%	\setlength\itemsep{1 in}
%	\item How much force would it take to lift a \SI{100}{kg} box straight up at constant speed a distance of 5 meters.
%	\item How much force would it take to push a box up an inclined plane that was \SI{20}{\meter} long up to the same height? (again at constant speed)
%	\item What is the ratio of the two forces in these two cases? Which would you rather do? This ratio is known as \emph{mechanical advantage}
%	\item What is the ratio of the displacement in these two cases?
%	\item How can you use this idea to quickly figure out the force it would take you to push the box up a \SI{100}{\meter} ramp that goes up the same height?
%\end{itemize}\hugeskip



%\item Let's talk about springs. \emph{Hooke's Law} is the relationship between the force exerted by a spring and amount the spring has been stretched or compressed. The amount a spring has been stretched or compressed is the displacement of the end of the spring. Hooke's law says the magnitude of the force exerted is directly proportional to the displacement of the end of the spring. Write this in terms of a proportionality statement and again as an equation with a \emph{constant of proportionality}. What are the units of the constant of proportionality?\bigskip
%
%Draw a qualitative plot of the magnitude of this force vs displacement. What is the slope of this graph?\hugeskip

\item Consider the graph below of displacement under condition of \emph{not constant} acceleration. What would a graph of velocity vs. time look like for this case? What about acceleration vs time. 

%\includegraphics[scale=.62]{non-constant-acceleration.pgf}\hugeskip
\input{non-constant-acceleration.pgf}



%\item 
%A follow up to the previous question. The \emph{momentum} of an object is defined as the product of an object's mass and its velocity. In this case the momentum of this object changes because a force external to the object was exerted on it. What is the momentum of the object initially? What is the momentum of the object at the end? It is said that the change in momentum of an object is equal to the \emph{impulse}, where \emph{impulse} is defined as the product of constant force and the amount of time the force is acting. Is that the case here? How could you use this to quickly find the final velocity of an object if the same force were pushing it for \SI{100}{\second}?
%\hugeskip

%\item 
%If I launch a cannonball from ground level with a speed of \SI{1000}{m/s}, then there are a variety of angles to choose from. If I fire it at \ang{90}, then how long will it take to come down? How high will it go? How far will it go horizontally? What about \ang{80}? What about \ang{70}? Let's just make a table...
%
%\begin{tabular}{rrrr}
%	\toprule
%	Angle & Time & Max Vertical Height & Horizontal Distance \\
%	\midrule
%	\midrule
%	\ang{90} \\
%	\ang{80} \\
%	\ang{70} \\
%	\ang{60} \\
%	\ang{50} \\
%	\ang{40} \\
%	\ang{30} \\
%	\ang{20} \\
%	\ang{10} \\
%	\ang{0} \\
%	\bottomrule
%\end{tabular}

\item If a soccer ball with a radius of \SI{10}{\centi\meter} is rolls along the ground without slipping at \SI{5}{m/s}, then how many revolutions does it roll through in \SI{10}{\second} and what distance has a point on the edge of the ball traveled?
Some starters:
\begin{itemize}
	\item How fast is it \emph{spinning}? By that we mean \emph{angular speed}.
	\item How many radians does the ball rotate through in this time? What is that in revolutions?
	\item How far does it roll in this time? Is this the same distance as the distance of a point on the edge of the ball? Why or why not?
\end{itemize}
\hugeskip

\item
Following up on the previous problem, how many seconds does it take for the ball to complete one revolution? This amount of time is referred to as the \emph{period} of its rotation, and this is a similar characteristic time for the motion of the ball as the \emph{period of a pendulum} was in the first lab.\hugeskip

\item 
Another follow up. How many revolutions does the ball travel through \emph{per second}? You could convert this from angular speed $\omega$ that you would have calculated in the first instance of this problem, but if all you new was the period of the ball's rotation, how could you calculate it from there? \emph{(Hint: what is the difference between revolutions per second and seconds per revolution?)} This quantity of revolutions per unit of time is sometimes called \emph{frequency}.\hugeskip

\item 
Suppose a satellite is in orbit around a distant planet. You observe the the satellite to be \SI{5000}{\kilo\meter} from the center of the planet, and rotating the planet once every 2 days. What is the mass of the planet you have discovered? What is the period of the satellites motion? What is its frequency? How fast is the satellite moving around the planet? What is the angular speed?
\bigskip

\item
A car travels at 17 m/s without skidding around a level curved road with a radius of 35m. What is the coefficient of static friction between the tires and the road if this speed is the fastest it can go around this curve without skidding?


\end{enumerate}
